\documentclass[UTF8,12pt,a4paper]{ctexart}
\usepackage{geometry}
\geometry{left=2.5cm,right=2.5cm,top=2.4cm,bottom=2.4cm}
\usepackage{amsmath,amssymb,bm}
\usepackage{graphicx}
\usepackage{booktabs,tabularx,array,multirow}
\usepackage{float}
\usepackage{caption}
\setlength{\parskip}{0.35em}
\setlength{\parindent}{2em}
\emergencystretch=1em
\captionsetup{font=small,labelfont=bf}
\newcommand{\pp}{\,\mathrm{pp}}
\title{面向 SAR 图像的可靠性引导 Mamba 状态传播方法\\
\large Reliability-Guided Selective State Propagation in Mamba for SAR Imagery}
\author{}
\date{论文工作稿：2026 年 10 月 6 日}
\begin{document}
\maketitle

\begin{abstract}
视觉状态空间模型以输入依赖的选择性状态更新实现近线性复杂度的长程依赖建模，但其更新参数由特征内容隐式生成，并未显式利用合成孔径雷达（SAR）图像中局部响应的结构线索——散斑起伏、背景杂波与孤立强散射可使不同结构意义的局部响应呈现相近的特征表现。本文提出面向 SAR 图像的可靠性引导选择性状态传播方法（Reliability-Guided Selective State Propagation, RGSP）：在特征进入选择性扫描之前，局部响应修正（Local Response Correction, LRC）依据局部趋势与残差调整待传播特征；局部结构引导模块（Local Structure Guidance Module, LSGM）进一步融合局部趋势与二维邻域变化生成位置相关的引导信号，显式条件化状态更新的写入强度与转移步长。在 SSDD 官方测试划分上，本文完成了五种结构变体、三个随机种子共 15 次从头训练的受控消融：局部修正与结构引导分别独立使用时，AP$_{50:95}$ 相对基线的配对平均增益为 $+1.95$ 与 $+1.09$ 个百分点，其中结构引导变体的召回提升最为明显（$+2.15$ 个百分点）。本文同时报告并分析了两组件朴素组合的负交互与 token 级引导不优于图像级对照的结果，为 SAR 状态空间建模的设计空间提供了可复现的组件级证据；分类、变化检测与分割任务上的诊断结果进一步界定了该设计的适用边界。
\end{abstract}
\noindent\textbf{关键词：}合成孔径雷达；视觉状态空间模型；Mamba；选择性状态更新；局部响应修正；目标检测

\section{引言}\label{sec:intro}

合成孔径雷达（Synthetic Aperture Radar, SAR）具有全天时、全天候成像能力以及较强的环境适应性，已广泛应用于海面舰船监测、机场目标识别、地表观测和灾害变化分析等场景。与自然光学图像不同，SAR 图像由目标散射特性、成像几何和相干成像过程共同决定，其局部响应通常伴随明显的散斑、背景杂波以及离散强散射现象。在复杂场景中，较强的局部响应既可能来自具有判别意义的目标散射中心，也可能由背景杂波或孤立强散射点产生；与此同时，部分响应较弱的区域仍可能包含目标边缘或细粒度结构信息。因此，SAR 局部响应的强弱并不能直接反映其所包含的任务相关结构信息。相比仅依据单个位置的响应进行特征建模，结合邻域关系和局部结构判断当前信息在目标表征中的作用，对于 SAR 视觉建模尤为重要。

另一方面，SAR 目标的判别信息通常并非集中于单一空间位置，而是由多个空间分布的散射中心、边缘以及几何结构共同构成。因此，模型不仅需要辨别局部信息的结构意义，还需要在更大的空间范围内建立不同区域之间的联系，以形成更加完整的目标表征。卷积神经网络（Convolutional Neural Network, CNN）通过层级结构逐步聚合局部特征并扩大有效感受野\cite{resnet2015}；视觉 Transformer 则利用自注意力直接建立不同图像区域之间的长距离依赖\cite{vit2020}，但全局成对交互在高分辨率特征上通常伴随较高的计算和存储开销。近年来，Mamba 及其选择性状态空间模型（Selective State Space Model, SSM）通过输入依赖的状态更新机制，以随序列长度近线性增长的计算复杂度实现长距离信息建模\cite{mamba2023}。Vision Mamba（Vim）将状态空间建模扩展到视觉序列\cite{vim2024}，VMamba 进一步通过二维选择性扫描（SS2D）建立空间依赖\cite{vmamba2024}，Mamba YOLO 等工作则将状态空间建模引入层级目标检测框架\cite{mambayolo2024}。这些研究表明，Mamba 为视觉特征的长距离传播与高效空间建模提供了一条区别于卷积和全局自注意力的新路径。

随着视觉状态空间模型的发展，相关研究进一步从空间序列组织和特征交互等方面增强 Mamba 的视觉建模能力。RS-Mamba 面向大幅面遥感图像研究状态空间模型在密集预测中的应用\cite{rsmamba2024}，RSMTMamba 探索跨任务遥感特征学习\cite{rstmtmamba2025}，MaDiNet 等方法则将 Mamba 引入 SAR 目标检测框架\cite{madinet2025}；此外，部分工作还从成像统计建模与输入特征构造等层面注入 SAR 领域先验\cite{mambadip,rssnet}。总体而言，现有视觉与遥感 Mamba 方法主要围绕二维或多方向扫描、多尺度特征交互以及卷积与状态空间模型融合等方面展开，重点改善空间序列如何组织、长距离依赖如何建立以及不同层级特征如何聚合。这些研究推动了 Mamba 在遥感视觉中的应用，也为 SAR 图像的长程结构建模提供了良好基础。

尽管 Mamba 具有输入依赖的选择性状态更新机制，将其直接应用于 SAR 图像仍存在一个值得进一步研究的问题。Mamba 的部分状态更新参数由输入特征动态生成，使模型能够根据当前内容调节信息写入与状态演化；然而，SAR 图像中的局部响应容易受到散斑、背景杂波和孤立强散射等因素影响，使具有不同结构意义的响应在局部特征上呈现出较强歧义。例如，显著响应既可能对应稳定的目标散射结构，也可能表现为缺少邻域结构支持的孤立响应，而部分相对较弱的位置仍可能包含连续边缘或细粒度结构。在这种情况下，内容相关的选择性更新可能受到局部响应歧义的影响，而现有 Mamba 并未显式利用 SAR 局部结构关系来调节不同响应进入状态更新的方式与强度。更重要的是，一旦某些局部信息参与当前状态更新，其影响便会进入递推状态，并沿扫描过程影响后续位置的特征表示。因此，对于局部散射结构复杂的 SAR 图像，一个关键问题在于：\emph{如何在信息进入递推状态之前利用局部结构改善其表示，并在后续状态更新中进一步依据结构差异调节不同局部信息的传播？}

基于这一观察，本文进一步关注选择性状态传播中的\emph{传播内容与状态更新}。对于 SAR 图像而言，局部响应是否值得参与传播，不应仅由其当前特征的显著程度决定，还应结合其与周围位置之间的局部趋势和邻域变化进行判断。为此，一方面，可以利用当前位置与局部趋势之间的关系调整进入状态传播的特征表示，使结构关系不同的局部响应在进入递推过程之前得到差异化处理；另一方面，还需要将局部结构信息进一步引入状态更新，使不同空间位置在当前信息写入和历史状态演化过程中产生差异化作用。本文将局部响应在多大程度上反映稳定、任务相关的结构（而非瞬时的散斑起伏或孤立强散射响应）概括为\textbf{局部结构可靠性}。该性质难以从单个位置的响应强度直接判断，但可以从局部趋势与邻域变化等结构关系中获取线索；本文以这些线索为出发点，改进 Mamba 的选择性状态传播。

基于上述分析，我们提出一种面向 SAR 图像的\textbf{可靠性引导选择性状态传播方法（Reliability-Guided Selective State Propagation, RGSP）}（图\ref{fig:motivation}）。RGSP 显式利用局部趋势和邻域变化刻画当前响应与周围结构之间的关系，并从特征进入扫描与状态更新两个层面将这些结构信息引入选择性递推过程。具体而言，RGSP 首先根据局部趋势及当前响应相对于该趋势的偏差，对进入选择性扫描的中间特征进行可学习的调整，使不同局部结构关系能够反映在待传播的特征表示中；随后进一步融合局部趋势与二维邻域变化生成位置相关的结构引导信息，并利用该信息调节选择性状态更新中的输入写入和状态演化。由此，RGSP 并非对局部响应进行简单的噪声/有效特征二分类，而是利用显式局部结构信息进一步区分不同响应的传播作用，使 Mamba 的状态传播由单纯的内容相关选择进一步扩展为局部结构引导的选择性传播。

具体而言，RGSP 首先设计\textbf{局部响应修正（Local Response Correction, LRC）}。LRC 在网络中间特征上沿扫描序列学习局部趋势，并利用当前响应相对于局部趋势的偏差构造特征残差，对进入选择性扫描之前的特征进行可学习的调整。该过程并不以恢复平滑的 SAR 图像为目标，而是利用局部关系重新组织进入状态空间建模的特征表示，使当前响应的建模不再完全依赖单个位置本身的响应强弱。

在此基础上，本文进一步设计\textbf{局部结构引导模块（Local Structure Guidance Module, LSGM）}。LSGM 融合局部趋势信息与二维邻域变化，生成位置相关的结构引导信息，用于刻画不同空间位置在当前任务中的相对结构作用。该引导信息进一步被引入选择性状态更新过程，并分别作用于输入驱动和离散化步长：输入驱动调节当前特征参与状态写入的强度，动态步长进一步参与历史状态演化与当前信息写入。由此，RGSP 分别从“进入状态传播的信息”和“信息参与状态更新的程度”两个层面对 Mamba 的选择性递推进行调节，使不同空间位置能够根据其局部结构特征以差异化方式参与后续状态传播。

通过上述两个层面的设计，RGSP 将 SAR 局部结构线索从特征表示进一步引入选择性状态更新参数的形成过程，使局部结构分析直接影响信息写入与状态演化。

为评估 RGSP 在不同 SAR 视觉任务中的有效性与适用性，本文在目标检测、图像分类、语义分割和变化检测任务上开展实验，并通过消融实验分析不同组成部分的实际作用。此外，我们进一步与代表性方法进行比较，从任务性能和计算代价两个方面评估所提出方法。

\begin{figure}[tbp]
\centering
\includegraphics[width=\linewidth]{figures/fig1_motivation.pdf}
\caption{\textbf{RGSP 的研究动机。}（a）以来自 SSDD 训练集的一幅图像为例，SAR 局部响应的强度与任务相关结构意义并不一致：标注舰船（①）具有稳定的结构；未标注的强散射响应（②）强度与之相当，但不承载任务相关的目标结构；弱响应区域（③）同样可能包含结构线索。（b）通用选择性状态更新的参数仅由特征内容生成，对上述响应不作显式区分；（c）RGSP 先经局部响应修正（LRC）调整待传播特征，再由局部结构引导模块（LSGM）生成位置相关的引导信号，条件化选择性状态更新的写入强度与转移步长（机制为设计示意，具体效果由受控实验检验）。}
\label{fig:motivation}
\end{figure}

本文的主要贡献如下：

\begin{enumerate}

\item 我们提出了一种面向 SAR 图像的可靠性引导选择性状态传播方法 RGSP。针对 SAR 局部响应具有较强歧义，而 Mamba 的选择性状态更新缺少针对局部结构关系的显式传播调节这一问题，RGSP 将局部结构信息引入递推状态建模，使不同响应能够依据其结构关系以差异化方式参与状态传播。

\item 我们设计了局部响应修正（Local Response Correction, LRC）和局部结构引导模块（Local Structure Guidance Module, LSGM）。LRC 利用局部趋势及其与当前响应之间的偏差调整待传播特征，LSGM 进一步融合局部趋势与二维邻域变化生成位置相关的结构引导信息，并据此调节选择性状态更新中的输入写入和状态演化，从而建立局部结构分析与状态传播之间的直接联系。

\item 我们在 SAR 目标检测、图像分类、语义分割和变化检测等不同视觉任务中对 RGSP 进行了系统评估，并结合模块消融、代表性方法比较以及计算效率分析，考察各组成设计的实际作用及所提出方法在不同 SAR 任务中的适用性。

\end{enumerate}


\section{相关工作}
\subsection{SAR 图像去斑与目标检测}
SAR 去斑研究通过统计滤波或学习式映射减轻散斑影响；目标检测研究则直接从强度图像学习目标和背景判别特征。去斑与检测之间存在任务权衡：降低局部起伏可能压制部分噪声，也可能削弱小目标和边缘。Speckle-Free SAR ship detection 等工作已经研究了去斑与检测之间的关系\cite{specklefree}。近期 RSS-Net 将 Mamba 的二维选择性扫描与 CNN 局部特征结合用于 SAR 去斑；MambaDIP 则将 Mamba 放入深度图像先验网络，并使用 Gamma 似然处理 SAR 复原\cite{mambadip,rssnet}。这些工作说明 Mamba 已进入 SAR 去斑研究，也强调长程建模需要与局部细节和成像统计共同考虑。本文不把局部标准差解释为物理散斑估计，也不声称本文模块具有保证保边的滤波性质；它是状态更新所用的特征线索。

SSDD 是公开 SAR 舰船检测数据集之一\cite{ssdd2017,ssdd2021}。

SAR-ShipDataset 与 SAR-AIRcraft1.0 则涵盖不同成像场景和目标类别\cite{sardataset,saraircraft}。这些数据集上的数字只有在划分、初始化、训练预算和评估口径一致时才能进行直接比较。本文因此将 SSDD 三种子受控消融作为当前主证据，将其他数据集的历史数字单独列为不可直接排名的记录。

\subsection{视觉状态空间模型与遥感 Mamba}
Mamba 通过选择性状态空间模型实现输入相关的状态更新\cite{mamba2023}；Vision Mamba 将该机制扩展到视觉表征\cite{vim2024}。遥感领域已有 RS-Mamba、跨任务 Mamba、Mamba-YOLO 等工作\cite{rsmamba2024,rstmtmamba2025,mambayolo2024}。针对 SAR，也已有 MaDiNet、MCMN、RSS-Net 和 MambaDIP 等直接包含 Mamba 的研究，以及 SAR-NanoShipNet、SFD-Net 等近期目标检测方法\cite{madinet2025,mcmn,rssnet,mambadip,sarnano,sfdnet}。因此，本文的定位不能建立在“首次将 Mamba 用于遥感/SAR”这一宽泛表述上；更具体的问题是，现有 Mamba-SAR 研究已经从 Gamma 统计先验、Doppler/散射输入、多尺度融合、复原损失等层面注入任务知识，而当前稿件研究的是特征进入 selective scan 后对扫描驱动和步长的调节，并需要通过对照实验判断其实际价值。

\subsection{输入依赖更新与 SAR 线索注入}
选择性状态空间模型由输入相关参数调节状态更新；视觉 Mamba 和遥感 Mamba 工作进一步研究了双向、全向或跨任务扫描路径\cite{mamba2023,vim2024,rsmamba2024,rstmtmamba2025}。SAR 领域的已有工作从不同层级注入先验：MaDiNet 将 Gamma 分布建模与 MambaSAR 用于目标检测；MCMN 在 PolSAR 船舶检测中构造 Doppler/散射输入特征，并采用 Mamba-CNN 局部与全局分支；RSS-Net 将二维选择性扫描与 CNN 局部特征结合用于 SAR 去斑；MambaDIP 则将 Mamba 放入深度图像先验 U-Net，并用 Gamma 似然处理 SAR 去斑与增强\cite{madinet2025,mcmn,rssnet,mambadip}。近期非 Mamba 检测器如 SAR-NanoShipNet 和 SFD-Net 则围绕小目标边界、多尺度和干扰特征分解设计\cite{sarnano,sfdnet}。这类研究已覆盖 SAR 检测与复原，因而本文不以“尚无 SAR-Mamba”为缺口。

表\ref{tab:related_position}按先验进入模型的位置归纳代表性方法。本文考察的是特征进入扫描后、状态更新参数形成时的调节位置；其差别是设计定位，不代表优于已发表方法。尤其是，当前 SSDD 对照仍含有 block 结构差异，故实验只能评估所实现变体的整体效果，尚不能把差异完全归因于单一公式项。

\begin{table}[H]\centering\scriptsize
\caption{代表性 SAR 状态空间、去斑及目标检测方法与本文研究位置。该表为相关工作抽样，不是穷尽性综述。}
\label{tab:related_position}
\begin{tabularx}{\linewidth}{p{3.2cm}p{2.25cm}p{3.7cm}X}
\toprule
方法 & 任务范围 & 主要先验或结构位置 & 与本文的关系\\
\midrule
Speckle-Free ship detection\cite{specklefree} & SAR 舰船检测 & 将去斑处理与检测流程结合 & 说明去斑和检测的任务权衡；不是状态空间更新方法\\
RS-Mamba\cite{rsmamba2024} & 大幅面遥感密集预测 & 全向选择性扫描以聚合多方向上下文 & 关注扫描方向与长程上下文，不专门建模 SAR 局部统计\\
MaDiNet\cite{madinet2025} & SAR 多类目标检测 & Gamma 分布建模与 MambaSAR 特征提取 & 与本文最直接相关；其先验涉及成像统计/特征建模，本文研究 scan 驱动和步长的调节\\
MCMN\cite{mcmn} & PolSAR 船舶检测 & Doppler/散射输入特征及 Mamba-CNN 局部与全局分支 & 在输入与网络特征层注入物理线索；本文的重点是状态更新路径\\
MambaDIP\cite{mambadip} & SAR 去斑与增强 & 深度图像先验 U-Net、Mamba 层与 Gamma 似然 & 属于图像复原；不能作为目标检测结果的直接基线\\
RSS-Net\cite{rssnet} & SAR 去斑 & VSSM 四向扫描与 CNN 残差特征块 & 与本文共享 SAR 状态空间建模背景；任务目标是复原而非检测\\
SAR-NanoShipNet\cite{sarnano} & SAR 小目标舰船检测 & 可变形卷积、边界注意与尺度适配 & 提供近期 SAR 检测设计参照，不使用 Mamba\\
SFD-Net\cite{sfdnet} & SAR 多尺度舰船检测 & 可学习特征解耦与 Star-Fusion FPN & 处理目标/干扰特征分离；不是状态空间更新方法\\
RSMTMamba\cite{rstmtmamba2025} & 遥感多任务预测 & 共享编码与跨任务 Mamba 特征交互 & 支持跨任务设计背景，但不等同于 SAR 噪声感知更新\\
本文 & SAR 视觉状态更新 & 局部残差修正；重要性调节 scan drive 与 $\Delta$ & 当前 SSDD 证据显示单模块正向均值、组合负交互；不作普遍领先主张\\
\bottomrule
\end{tabularx}
\end{table}

\subsection{本文定位与可检验假设}
我们将局部趋势、残差和邻域起伏作为 SAR 特征更新的候选线索，提出以下假设：局部残差修正可能减少无关高频响应对扫描输入的影响；重要性调制可能对疑似目标相关特征分配不同的写入强度与状态步长。它们是方法动机而非实证结论。本文通过单模块、组合和粒度对照检验这些设计，并保留反向或不稳定的结果。
\section{方法}
\label{sec:method}

\subsection{总体网络与 RGSP 块}
\label{sec:method_overview}

本文在视觉状态空间网络中引入可靠性引导选择性状态传播机制（Reliability-Guided Selective State Propagation，RGSP），利用显式局部结构线索调节选择性状态更新。如图~\ref{fig:overview}(a) 所示，任务网络由输入嵌入、分层骨干、特征聚合或解码模块以及任务头组成。SAR 图像经过输入嵌入与骨干网络提取层级特征，再由任务头生成预测。RGSP 集成于骨干的状态空间块中，调整中间特征的递归建模过程。

图中以检测配置为例，将 RGSP 块部署于第二个骨干阶段的 P3/8 特征层，其输出继续传入后续骨干阶段、特征聚合模块和检测头。分类、检测和分割标签表示可独立配置的任务网络，而非同一模型的并行输出；不同任务的实际接入方式在实验部分说明。本节以该骨干内部替换块为对象，介绍其块级结构及内部状态建模过程。

\begin{figure}[t]
\centering
\includegraphics[width=\linewidth]{figures/fig2_method.pdf}
\caption{整体任务网络与 RGSP 的内部信息流。（a）以检测网络为例，高亮 P3/8 特征层中的 RGSP 块，并通过虚线连接下方放大图。任务头标签表示不同任务的独立配置。（b）依次展示局部趋势与残差修正（LRC）、局部结构引导（LSGM）和结构条件化状态更新。深色实线为特征流，橙色虚线为结构引导路径。图中聚合展示扫描后的旁路、门控和投影；完整块的残差与前馈结构见式~\eqref{eq:rgsp_block}。}
\label{fig:overview}
\end{figure}

\paragraph{块级结构。}
给定中间特征 $F\in\mathbb{R}^{b\times H\times W\times C}$，将其按行优先顺序展平为 $X\in\mathbb{R}^{b\times N\times C}$，其中 $N=HW$。RGSP 块由结构条件化状态建模分支与前馈分支组成，两者依次通过残差连接更新特征：
\begin{equation}
\begin{aligned}
U&=X+a_1\odot\mathcal{R}\bigl(\operatorname{LN}_1(X)\bigr),\\
X_{\mathrm{out}}&=U+a_2\odot\operatorname{FFN}\bigl(\operatorname{LN}_2(U)\bigr).
\end{aligned}
\label{eq:rgsp_block}
\end{equation}
其中，$\mathcal{R}$ 表示 RGSP 的状态建模分支，$\operatorname{FFN}$ 为带 GELU 激活的两层前馈网络，$a_1,a_2\in\mathbb{R}^{C}$ 为可学习的逐通道缩放参数。最后将 $X_{\mathrm{out}}$ 恢复为空间特征图，传入后续网络层。

\paragraph{分支内部流程。}
如图~\ref{fig:overview}(b) 所示，$\mathcal{R}$ 首先通过局部趋势与残差修正（Local Trend and Residual Correction，LRC）形成修正特征；随后由局部结构引导模块（Local Structure Guidance Module，LSGM）联合趋势与二维邻域变化生成结构引导信号；最后将该信号分别注入扫描输入与离散化步长，执行选择性递推并生成分支输出。以下记 $x=\operatorname{LN}_1(X)$，按上述数据流依次展开。可靠性是采用局部结构线索的设计动机，实际学习的引导信号由下游任务损失优化，用于控制状态更新。

\subsection{局部趋势与残差修正}
\label{sec:lrc}

RGSP 首先对归一化序列 $x$ 建立短程邻域描述。LRC 沿扫描序列执行核宽为 $3$ 的逐通道一维卷积，估计局部趋势：
\begin{equation}
\bar{x}=\mathcal{L}(x)=\operatorname{DWConv1D}_{k=3}(x).
\label{eq:local_trend}
\end{equation}
各通道独立聚合短程信息，卷积参数与任务网络联合学习。在此基础上，通过趋势残差调整扫描前的特征：
\begin{equation}
r=x-\bar{x},\qquad \tilde{x}=x-\lambda_s r.
\label{eq:residual}
\end{equation}
其中，$\lambda_s$ 为块级可学习标量，控制特征相对于局部趋势的偏离对修正结果的影响。趋势卷积与修正系数均参与训练，使邻域关系能够反映在后续内容表示中。

LRC 输出两路信息：$\bar{x}$ 为结构描述和趋势能量提供输入，$\tilde{x}$ 用于二维邻域统计与内容投影。该趋势算子沿行优先展平序列计算，图像行末与下一行行首会成为相邻 token；因此它描述与扫描顺序一致的一维短程关系，二维空间信息由 LSGM 补充。

\subsection{局部结构引导}
\label{sec:lsgm}

获得 $\bar{x}$ 与 $\tilde{x}$ 后，LSGM 将序列趋势和二维邻域变化联合编码为结构引导信号。首先将修正特征恢复为空间布局 $\widetilde{F}$，逐通道计算 $3\times3$ 邻域标准差：
\begin{equation}
S=\sqrt{\max\bigl(\mathcal{P}(\widetilde{F}^{\,2})-\mathcal{P}(\widetilde{F})^{2},0\bigr)+\epsilon}.
\label{eq:local_std}
\end{equation}
其中，平方运算逐元素执行，$\mathcal{P}$ 为局部均值池化，$\epsilon$ 为数值稳定常数；边界处仅使用有效邻域位置计算统计量。将 $S$ 按相同空间顺序展平为 $s$，与趋势特征通道拼接：
\begin{equation}
v=[\bar{x},s]\in\mathbb{R}^{b\times N\times2C}.
\label{eq:structural_descriptor}
\end{equation}
随后通过归一化与两层感知机生成引导信号：
\begin{equation}
I=\sigma\!\left(W_2\operatorname{GELU}\!\left(W_1\operatorname{LN}(v)+b_1\right)+b_2\right),
\quad I\in(0,1)^{b\times N\times1}.
\label{eq:importance}
\end{equation}
其中 $\sigma$ 为 Sigmoid 函数，$W_1,W_2,b_1,b_2$ 为可学习参数。LSGM 的隐藏宽度为 $\max(\lfloor C/4\rfloor,16)$，每个位置输出一个标量。趋势与局部变化通过任务训练联合解释，不预设局部变化程度与信息价值的单调关系。所得 $I$ 是结构条件信号，在后续操作中沿通道维广播，分别进入扫描输入调制和离散化步长条件化路径，如图中的两条橙色虚线所示。

\subsection{结构条件化选择性状态更新}
\label{sec:structure_conditioned_ssm}

LRC 与 LSGM 分别提供修正特征 $\tilde{x}$ 和结构引导信号 $I$。本节说明二者如何共同形成扫描输入与动态参数，并通过递归计算产生状态建模分支的输出。以下省略批维，以 $t$ 表示序列位置，以 $d_s$ 表示每个通道的状态维数。

\paragraph{内容投影。}
首先将修正特征投影为内容分支与输出门控分支：
\begin{equation}
[u_t,z_t]=W_{\mathrm{in}}\tilde{x}_t,\qquad u_t,z_t\in\mathbb{R}^{C}.
\label{eq:input_projection}
\end{equation}
内容分支进一步生成状态写入与读出参数：
\begin{equation}
B_t=W_Bu_t,\qquad C_t=W_Cu_t,\qquad B_t,C_t\in\mathbb{R}^{d_s}.
\label{eq:ssm_parameters}
\end{equation}
$B_t$ 和 $C_t$ 均由输入调制前的 $u_t$ 生成。结构引导 $I_t$ 显式接入下述两条调制路径，而不直接接入 $B_t,C_t$ 的投影。

\paragraph{扫描输入调制。}
由局部趋势定义逐位置能量：
\begin{equation}
E_t=\frac{\|\bar{x}_t\|_2}{\sqrt{C}}.
\label{eq:trend_energy}
\end{equation}
结合 $E_t$ 与 $I_t$ 构造两个输入调制因子：
\begin{align}
\gamma_t&=g_{\min}+(1-g_{\min})\sigma\!\left(-\alpha s_nE_t+\beta(I_t-0.5)\right),\\
g_{w,t}&=\max\!\left(0.1,1+\tanh(\rho)(I_t-0.5)\right),\\
u_{\mathrm{scan},t}&=\gamma_tg_{w,t}u_t.
\label{eq:modulation}
\end{align}
其中 $\alpha,\beta,\rho$ 为可学习标量，$s_n$ 为固定能量缩放系数，$g_{\min}$ 为 $\gamma_t$ 的下界。两个因子共同作用于同一内容分支，形成实际送入选择性扫描的输入驱动。图中的趋势能量支路与结构引导支路在此汇合。

\paragraph{离散化步长条件化。}
在内容相关步长的 logits 上加入结构引导项：
\begin{equation}
\Delta_t=\operatorname{softplus}\!\left(W_{\Delta}u_t+b_{\Delta}+\eta I_t\right),
\qquad\Delta_t\in\mathbb{R}_{>0}^{C},
\label{eq:delta_condition}
\end{equation}
其中 $\eta$ 为可学习标量，$I_t$ 沿通道维广播。该路径使递推中的离散化步长同时依赖内容特征与显式局部结构描述。

\paragraph{选择性递推与输出。}
状态矩阵采用 $A=-\exp(A_{\log})\in\mathbb{R}^{C\times d_s}$ 的参数化。令 $h_{t,c,j}$ 表示位置 $t$、通道 $c$ 和状态维 $j$ 上的隐藏状态，选择性递推为
\begin{equation}
h_{t,c,j}=\exp(\Delta_{t,c}A_{c,j})h_{t-1,c,j}+\Delta_{t,c}B_{t,j}u_{\mathrm{scan},t,c},
\quad h_{0,c,j}=0.
\label{eq:rgsp_state}
\end{equation}
由状态读出并加入修正特征旁路：
\begin{equation}
y_{t,c}=\sum_{j=1}^{d_s}C_{t,j}h_{t,c,j}+D_c\tilde{x}_{t,c},
\label{eq:ssm_readout}
\end{equation}
其中 $D\in\mathbb{R}^{C}$ 为可学习参数。随后执行输出门控和投影：
\begin{equation}
o_t=W_{\mathrm{out}}\bigl(y_t\odot\operatorname{SiLU}(z_t)\bigr).
\label{eq:rgsp_output}
\end{equation}
上述读出、旁路、门控及投影对应图中的 ``Skip + Gate + Projection''。$u_{\mathrm{scan},t}$ 调节当前内容的输入驱动，$\Delta_t$ 同时参与历史状态衰减与当前输入写入；由于 $A$ 的元素为负，增大步长会增强相应历史分量的衰减，其作用需结合输入项共同理解。

将各位置输出 $o_t$ 组成序列，即得到式~\eqref{eq:rgsp_block} 中状态建模分支 $\mathcal{R}$ 的输出。该序列经块级残差融合与前馈分支形成最终特征。由此，局部结构分析通过扫描输入和离散化步长两条路径参与状态递推，并将所得表示传递至后续网络层。


\section{实验设置}
\subsection{主数据集与划分}
主验证数据集为 SSDD。官方训练部分含 928 张图像，官方测试集含 232 张图像。对每个 seed，训练部分按冻结的 manifest 划分为 742 张训练图像和 186 张验证图像；同一 seed 下的五个变体共享这一划分，seed 0、1、2 之间则使用不同划分。因此，逐 seed 配对差控制同一划分，跨 seed 方差同时包含划分变化和随机初始化变化。测试集只用于最终评估，模型选择依赖验证集；没有使用 test-as-val。

分类、变化检测和分割作为探索性跨任务记录单独汇总。SAR-ShipDataset 与 SAR-AIRcraft1.0 虽有历史单次结果，但尚无 baseline、LRC-only、LSGM-only 和 full 四种变体在同一划分、scratch 初始化及相同预算下的三种子矩阵；它们不进入当前主比较。各探索任务的协议和数据覆盖见表\ref{tab:evidence_scope}及附录。

\subsection{统一训练协议与选模}
SSDD 主实验的五个变体均从头训练（\texttt{weights=none}），输入尺寸为 $640$，batch size 为 $16$，最多训练 150 epochs，patience 为 30。优化器为 AdamW，初始学习率 $2\times10^{-4}$、最终学习率比例 0.01、weight decay $5\times10^{-4}$、名义 batch size（nbs）64；使用 cosine learning-rate schedule，统一 warmup 3 epochs，新增模块 warmup 为 0。数据增强统一为 mosaic 0.3（最后 15 epochs 关闭）、degrees 180、translate 0.1、scale 0.3、vertical/horizontal flip 各 0.5，以及 HSV 参数 $(0,0,0.1)$；启用确定性设置，不使用 AMP。

实验在两张 RTX A6000 上执行，三个 seed 均使用相同型号 GPU；每个 seed 的五个变体在同一张卡上完成。运行环境记录为 Python 3.10.16、PyTorch 1.13.1+cu117、Ultralytics 8.2.29，本地 selective-scan CUDA kernel 使用 FP32。统一依据验证集 AP$_{50:95}$ 选择 best checkpoint，再对官方测试集评估；所有变体遵循相同选模和评估规则。逐 run 指标及权重 SHA-256 保存在随稿补充材料中。

图\ref{fig:protocol}概括了主实验的数据隔离、逐 seed 配对和选模流程。分类、变化检测、分割及历史运行不进入这条受控矩阵，单独按其各自协议报告。
\begin{figure}[H]\centering
\includegraphics[width=.98\linewidth]{figures/fig11_experiment_protocol.pdf}
\caption{SSDD 受控实验流程。五个变体在每个 seed 的同一训练/验证划分上训练；验证集负责选模，官方测试集只在最终评估阶段使用。}
\label{fig:protocol}
\end{figure}

\subsection{对比方法与受控变体}
五个主变体均使用相同的 Mamba-YOLO-T 主干、检测颈部和任务头。基线保留原始 VSSBlock；其余变体将 P3/8 阶段的 VSSBlock 替换为对应的 RGSP 块（该阶段经深度缩放后实际执行一个块），状态维数为 16。变体语义与实际开关列于表\ref{tab:variants}。该对照评估的是已实现块替换的整体效果：RGSP 块与原 VSSBlock 在归一化、扫描分支及 MLP ratio 上并不完全相同，本文不将差异归因于单一残差公式。

当前没有完成 YOLOv8n、YOLOv5s 等轻量检测器的同协议重实现，也没有完成固定 Lee 去斑预处理对照。已有文献报告值和历史运行记录分别标注来源与协议限制，不伪装成统一预算下的实测排行榜或预处理消融。
\begin{table}[H]\centering\scriptsize
\caption{SSDD 受控变体及其实现开关。}
\label{tab:variants}
\begin{tabularx}{\linewidth}{p{2.5cm}p{3.2cm}p{3.0cm}X}
\toprule
变体 & 替换位置/模块 & LRC 与能量门控 & LSGM 路径\\
\midrule
Baseline & 原始 VSSBlock & 关闭 & 关闭\\
LRC-only P3 & P3 的 RGSP 块 & 开启 & LSGM 关闭\\
LSGM-only P3 & P3 的 RGSP 块 & $\lambda_s=0,\alpha=0$ 且冻结 & token-wise；保留 $\gamma$/write/$\Delta$ 调制\\
Full P3 & P3 的 RGSP 块 & 开启 & token-wise；保留各调制路径\\
Global P3 & P3 的 RGSP 块（image-wise） & 开启 & image-wise；用于引导信息粒度对照\\
\bottomrule
\end{tabularx}
\end{table}

{\footnotesize 注：变体名与实现代码的对应关系为 LRC-only$\leftrightarrow$\texttt{ds\_only\_p3}、LSGM-only$\leftrightarrow$\texttt{pfim\_only\_p3}、Full$\leftrightarrow$\texttt{full\_p3}、Global$\leftrightarrow$\texttt{global\_p3}；LSGM 在实现代码中记为 PFIM。}

\subsection{指标、统计口径与效率测量}
检测报告 AP$_{50}$、AP$_{75}$、COCO 风格 AP$_{50:95}$、precision 和 recall。表中 mean$\pm$std 基于三个 seed；模块效果另报告同 seed 相对 baseline 的配对差。由于 $n=3$ 且 split 随 seed 变化，本文将其作为描述统计，不声称统计显著性。参数量和 GFLOPs 为记录的 seed 0 复杂度值；延迟和 FPS 虽保留在账本表中，但计时设备与复测元数据不完整，不作可靠速度排名。

\subsection{探索性跨任务与历史证据}
分类、变化检测和分割采用不同模型头、数据集接口或训练协议，且多为单 seed。它们在本文中被标记为探索性/部分可复现记录，不与 SSDD 主表合并，不支撑统一的“通用提升”结论。历史 SAR-ShipDataset、SAR-AIRcraft1.0、SADD 数字的训练划分或初始化口径不完全一致，也只在单独表中说明来源限制。旧 SSDD ex27 使用非 scratch 初始化且存在验证集口径差异，不作为主证据。

为避免把实验计划当作已完成结果，表\ref{tab:evidence_scope}汇总本文使用的证据层级；不同任务指标和协议不做跨行排名。
\begin{table}[H]\centering\scriptsize
\caption{任务范围、现有证据与本文中的用途。}
\label{tab:evidence_scope}
\setlength{\tabcolsep}{4.5pt}
\begin{tabularx}{.99\linewidth}{p{2.1cm}p{2.4cm}p{4.8cm}X}
\toprule
任务/数据集 & 角色 & 当前记录 & 证据边界\\
\midrule
SSDD 检测 & 主验证 & 官方 test；5 变体 $\times$ 3 seed；clean scratch & 当前唯一完整多种子受控矩阵\\
\shortstack[l]{SAR-Ship\\Dataset} 检测 & 计划主验证 & 有历史单次记录；正式公平三种子矩阵未完成 & 历史指标不得与 SSDD 合并排名\\
\shortstack[l]{SAR-AIRcraft\\1.0} 检测 & 计划主验证 & 有 patch 协议历史记录；正式公平三种子矩阵未完成 & 划分/训练协议需统一后复验\\
SAR-ACD、SAMPLE 分类 & 跨任务诊断 & 有单种子 scratch 与初始化/预热记录 & 不同训练轨道分开报告，不形成稳定性结论\\
Ottawa、Farmland、Yellow River、Sulzberger 变化检测 & 跨任务诊断 & 多种 batch/state dimension 的单种子记录 & 数据集间方向混合，批次之间不作重复样本统计\\
FUSAR-Map、\shortstack[l]{OpenEarthMap-\\SAR} 分割 & 探索边界 & 有受控单种子结果，mIoU 未改善；FUSAR road 类明显退化 & 不支撑通用性增益\\
\bottomrule
\end{tabularx}
\end{table}
\section{实验结果}
\subsection{SSDD 三种子主消融}
表\ref{tab:ssdd}报告三个 seed 的均值与标准差。LRC-only 在 AP$_{50:95}$ 上比 baseline 高 1.95 个百分点，LSGM-only 高 1.09 个百分点；full 高 0.88 个百分点，但低于 LRC-only 1.07 个百分点。global/image-wise 变体的均值比 full/token-wise 高 0.86 个百分点。AP$_{50}$ 与 AP$_{75}$ 指标呈现不同排序，故不能据单一指标概括“全面提升”。
\begin{table}[H]\centering\small
\caption{SSDD 五变体官方测试结果，三个 seed，mean$\pm$std（单位：\%）。}
\label{tab:ssdd}
\begin{tabular}{lrrrrr}
\toprule
Variant & AP$_{50}$ & AP$_{75}$ & AP$_{50:95}$ & Precision & Recall\\
\midrule
Baseline & $86.52\pm2.56$ & $34.05\pm5.25$ & $42.03\pm3.47$ & $89.55\pm1.83$ & $79.24\pm1.70$\\
LRC-only & $86.70\pm1.01$ & $37.58\pm3.78$ & $43.98\pm1.44$ & $89.89\pm0.96$ & $79.67\pm0.26$\\
LSGM-only & $88.23\pm1.61$ & $34.83\pm3.94$ & $43.12\pm1.70$ & $90.13\pm1.27$ & $81.39\pm2.20$\\
Full & $87.97\pm0.11$ & $34.10\pm3.12$ & $42.91\pm1.13$ & $91.91\pm1.49$ & $79.85\pm1.02$\\
Global/image-wise & $86.41\pm0.66$ & $36.87\pm3.89$ & $43.77\pm1.65$ & $91.76\pm0.85$ & $78.45\pm0.90$\\
\bottomrule
\end{tabular}
\end{table}
\begin{figure}[H]\centering
\includegraphics[width=.78\linewidth]{figures/fig3_ssdd_main.pdf}
\caption{SSDD 主指标的 mean$\pm$std。}
\label{fig:ssdd}
\end{figure}

\subsection{逐 seed 配对变化与模块交互}
表\ref{tab:paired}以百分点列出每个变体相对同 seed baseline 的 AP$_{50:95}$ 变化。LRC-only 和 LSGM-only 均在两个 seed 为正、一个 seed 为负；full 也仅两个 seed 的均值有一个明显正向种子，其余两次为负。LRC-only 的配对均值最高；受限于三个种子，本文以描述统计呈现上述差异。
\begin{table}[H]\centering\small
\caption{相对同 seed baseline 的 AP$_{50:95}$ 配对变化（百分点）。}
\label{tab:paired}
\begin{tabular}{lrrrr}
\toprule
Variant & Seed 0 & Seed 1 & Seed 2 & Mean$\pm$std\\
\midrule
LRC-only & $+2.90$ & $+5.26$ & $-2.30$ & $+1.95\pm3.87$\\
LSGM-only & $+2.13$ & $+3.16$ & $-2.00$ & $+1.09\pm2.73$\\
Full & $-0.35$ & $+5.93$ & $-2.94$ & $+0.88\pm4.57$\\
Global/image-wise & $+2.29$ & $+6.15$ & $-3.23$ & $+1.74\pm4.71$\\
\bottomrule
\end{tabular}
\end{table}
\begin{figure}[H]\centering
\includegraphics[width=.78\linewidth]{figures/fig4_ssdd_paired_delta.pdf}
\caption{各 seed 的配对 AP$_{50:95}$ 变化。图中保留负向 seed，避免只展示平均数掩盖波动。}
\label{fig:paired}
\end{figure}
full 相对 LRC-only 下降 1.07 个百分点，说明在当前实现和训练设定下，简单叠加两种改动没有产生互补增益。该结果与 LSGM 单独版本的正向均值并不矛盾：两个模块单独有效不保证组合有效。可能解释包括调制冗余、优化交互或 block 结构差异，但当前没有直接机制实验可区分这些解释。token-wise full 低于 global/image-wise 对照也表明，更细粒度空间预测在当前数据量和实现下没有展示优势。

\subsection{参数量与计算记录}
\begin{table}[H]\centering\small
\caption{SSDD 复杂度账本记录（seed 0）。速度口径元数据不完整，表中数值仅作记录，不作为已验证性能优势。}
\label{tab:efficiency}
\begin{tabular}{lrrrr}
\toprule
Variant & Parameters & GFLOPs & FPS & Latency (ms)\\
\midrule
Baseline & 5,985,547 & 13.61 & 151.1 & 6.62\\
LRC-only & 5,918,104 & 13.12 & 166.7 & 6.00\\
LSGM-only & 5,920,441 & 13.15 & 163.8 & 6.11\\
Full & 5,920,441 & 13.15 & 164.2 & 6.09\\
Global/image-wise & 5,920,441 & 13.15 & 163.2 & 6.13\\
\bottomrule
\end{tabular}
\end{table}
\begin{figure}[H]\centering
\includegraphics[width=.76\linewidth]{figures/fig5_efficiency.pdf}
\caption{SSDD AP$_{50:95}$ 与记录 GFLOPs 的关系。各变体参数量的精确值列于表\ref{tab:efficiency}；速度未绘入图中，因为原始计时元数据不完整。}
\label{fig:efficiency}
\end{figure}
参数量与 GFLOPs 记录显示，改造变体没有增加整体模型容量，部分指标略低于 baseline。不过 baseline 与 LRC-only 的 MLP ratio 和块内部结构不同；复杂度差异不是单独模块开销的纯估计。投稿前需用同硬件、同精度、batch=1 和同一 profiler 对完整检测器重新测量，并报告预热轮次、重复次数、同步方式和硬件型号。
\subsection{跨任务探索性结果与负面证据}
\paragraph{SAR 分类。}
分类实验存在两个不能合并的单种子轨道。scratch 轨道在 SAR-ACD 与 SAMPLE 上分别低于对应 baseline；另一条使用 VMamba-Tiny 初始化的 SAR-ACD 轨道记录了 P3 warmup 后的正向结果，但 warmup 改变训练日程。表\ref{tab:cls}将它们分开，避免把初始化与训练日程差异写成模块的稳定因果效应。旧六分类结果与当前五分类口径不兼容。
\begin{table}[H]\centering\scriptsize
\caption{分类跨任务单种子记录（Accuracy / Macro-F1，单位：\%）。不同轨道不可相互排名。}
\label{tab:cls}
\begin{tabularx}{\linewidth}{p{1.6cm}p{2.6cm}p{2.6cm}p{2.6cm}p{1.5cm}X}
\toprule
轨道 & 数据集/设置 & Baseline & P3 & $\Delta$ Accuracy & 说明\\
\midrule
Scratch & SAR-ACD & 94.882 / 94.815 & 89.370 / 89.181 & $-5.512$ & 224，BS64，至多 200e\\
Scratch & SAMPLE & 98.519 / 98.454 & 97.037 / 96.441 & $-1.482$ & real-only 选样口径\\
VMamba init & \shortstack[l]{SAR-ACD\\no warmup} & 89.764 / 89.585 & 91.339 / 91.275 & $+1.575$ & 新模块随机初始化\\
VMamba init & \shortstack[l]{SAR-ACD\\warmup 20} & 89.764 / 89.585 & 96.457 / 96.410 & $+6.693$ & P3 专属预热\\
VMamba init & \shortstack[l]{SAR-ACD\\warmup 50} & 89.764 / 89.585 & 96.850 / 96.830 & $+7.086$ & P3 专属预热\\
\bottomrule
\end{tabularx}
\end{table}
\begin{figure}[H]\centering
\includegraphics[width=.72\linewidth]{figures/fig7_sar_acd.pdf}
\caption{分类 seed 0 单次记录：左侧为 scratch 的 SAR-ACD/SAMPLE 配对，右侧为 VMamba-Tiny 初始化的 SAR-ACD 轨道。右侧 warmup 改变了 P3 训练日程，不能作为同预算的模块比较。}
\label{fig:acd}
\end{figure}

\paragraph{变化检测。}
变化检测结果均为 seed 0。P3 $d=16$ 的 BS16/BS96 记录与后续 $d=32$, BS16 批次是独立实验，不作重复样本统计。Farmland 在这些记录中均有较大正差；Yellow River 为负；Ottawa 的方向和幅度较小；Sulzberger 记录为负。表\ref{tab:cd}列出 $d=16$ 记录，表\ref{tab:cd32}单列后来完成的 $d=32$ 记录。它们支持结果依赖数据集与配置的描述，不构成多种子或普遍提升证据。
\begin{table}[H]\centering\scriptsize
\caption{变化检测历史结果（seed 0；每行的 batch 条件单独比较；单位：\%）。}
\label{tab:cd}
\begin{tabular}{llrrrrrr}
\toprule
Dataset & Method & Batch & Precision & Recall & F1 & IoU & OA\\
\midrule
Ottawa & Baseline & 16 & 90.259 & 94.000 & 92.091 & 85.342 & 98.351\\
Ottawa & P3 & 16 & 88.777 & 96.003 & 92.249 & 85.612 & 98.352\\
Ottawa & Baseline & 96 & 87.953 & 95.147 & 91.409 & 84.177 & 98.173\\
Ottawa & P3 & 96 & 89.841 & 94.346 & 92.039 & 85.251 & 98.333\\
Yellow River & Baseline & 16 & 91.205 & 85.702 & 88.368 & 79.160 & 97.817\\
Yellow River & P3 & 16 & 91.422 & 82.787 & 86.891 & 76.820 & 97.584\\
Yellow River & Baseline & 96 & 89.377 & 78.320 & 83.484 & 71.650 & 97.002\\
Yellow River & P3 & 96 & 91.218 & 71.146 & 79.941 & 66.585 & 96.546\\
Farmland & Baseline & 16 & 77.704 & 13.887 & 23.564 & 13.355 & 92.778\\
Farmland & P3 & 16 & 93.259 & 49.213 & 64.428 & 47.523 & 95.644\\
Farmland & Baseline & 96 & 80.623 & 12.532 & 21.692 & 12.165 & 92.747\\
Farmland & P3 & 96 & 91.747 & 50.719 & 65.325 & 48.506 & 95.684\\
Sulzberger & Baseline & 16 & 75.922 & 69.540 & 72.591 & 56.974 & 90.461\\
Sulzberger & P3 & 16 & 51.656 & 92.066 & 66.180 & 49.454 & 82.908\\
\bottomrule
\end{tabular}
\end{table}
\begin{figure}[H]\centering
\includegraphics[width=.79\linewidth]{figures/fig6_change_detection.pdf}
\caption{变化检测 P3 相对 baseline 的 F1 差（百分点）。左图展示 Ottawa、Yellow River、Sulzberger 的非 Farmland 记录，右图单独展示数值范围较大的 Farmland；颜色对应独立的 batch/state-dimension 设置。均为单 seed。}
\label{fig:cd}
\end{figure}
后续完成的 $d=32$, BS16 同协议数据集内配对结果为：Ottawa F1 92.758/92.110（$\Delta=-0.648$ pp），Yellow River 88.079/87.689（$\Delta=-0.390$ pp），Farmland 26.050/73.104（$\Delta=+47.054$ pp）；顺序均为 baseline/P3。该批次仍是单种子，且状态维数不同于表\ref{tab:cd}。
\begin{table}[H]\centering\small
\caption{$d=32$, BS16 变化检测结果（seed 0；单位：\%）。}
\label{tab:cd32}
\begin{tabular}{lrrrr}
\toprule
Dataset & Baseline F1 & P3 F1 & $\Delta$ F1 (pp) & P3 Params / GFLOPs\\
\midrule
Ottawa & 92.758 & 92.110 & $-0.648$ & 2.144M / 1.137\\
Yellow River & 88.079 & 87.689 & $-0.390$ & 2.144M / 1.137\\
Farmland & 26.050 & 73.104 & $+47.054$ & 2.144M / 1.137\\
\bottomrule
\end{tabular}
\end{table}

\paragraph{分割。}
FUSAR-Map 四类 UPerNet-lite 记录中\cite{fusar}，baseline 的 mIoU/FWIoU/OA 为 44.297/65.382/77.105\%，P3 为 43.561/65.083/76.770\%。P3 的 road IoU 从 4.782\% 降至 1.091\%，虽然 building 和 water IoU 略升。这是对任务适用边界的重要负面结果。

OpenEarthMap SAR-only 记录中，baseline mIoU 为 36.551\%；P3 无 warmup、warmup 20 和 50 的 mIoU 分别为 35.898\%、35.829\% 和 35.535\%，均低于 baseline。warmup 20 的 OA 52.499\% 略高于 baseline 52.316\%，但其 mPA 和 Dice 更低。该实验为 seed 0 部分可复现记录，不能把单个指标的微小变化描述成总体改善。
\begin{table}[H]\centering\small
\caption{历史分割实验摘要（单位：\%；单 seed / 部分可复现）。}
\label{tab:seg}
\begin{tabular}{llrrrrr}
\toprule
Dataset & Variant & mIoU & FWIoU & mPA & Dice & OA\\
\midrule
FUSAR-Map & Baseline & 44.297 & 65.382 & --- & --- & 77.105\\
FUSAR-Map & P3 & 43.561 & 65.083 & --- & --- & 76.770\\
OpenEarthMap-SAR & Baseline & 36.551 & 36.120 & 58.090 & 50.811 & 52.316\\
OpenEarthMap-SAR & P3, no warmup & 35.898 & 35.220 & 54.914 & 50.163 & 51.764\\
OpenEarthMap-SAR & P3, warmup 20 & 35.829 & 36.447 & 49.055 & 49.409 & 52.499\\
OpenEarthMap-SAR & P3, warmup 50 & 35.535 & 36.104 & 49.563 & 49.099 & 52.383\\
\bottomrule
\end{tabular}
\end{table}
\begin{figure}[H]\centering
\includegraphics[width=.83\linewidth]{figures/fig8_segmentation_limits.pdf}
\caption{分割历史结果中 P3 与 baseline 的主要指标差异。两个任务均未显示 mIoU 增益。}
\label{fig:seg}
\end{figure}
\paragraph{历史跨数据集检测。}
表\ref{tab:histdet}所列数字来自历史报告，协议存在训练轮数、初始化或划分方面的不一致，故仅作为需要复核的账本条目。它们不能与 SSDD 当前三种子结果合并排序，也不支撑跨数据集泛化结论。
\begin{table}[H]\centering\small
\caption{历史检测记录摘要（单位：\%；协议非统一，不可直接排名）。}
\label{tab:histdet}
\begin{tabular}{llrrl}
\toprule
Dataset / split & Method & AP$_{50}$ & AP$_{50:95}$ & Evidence note\\
\midrule
SAR-AIRcraft patch val & Baseline & 87.751 & 69.904 & 历史单次记录\\
SAR-AIRcraft patch val & P3 & 88.825 & 70.513 & 历史单次记录\\
SAR-AIRcraft patch test & Baseline & 70.402 & 54.513 & 历史单次记录\\
SAR-AIRcraft patch test & P3 & 70.822 & 55.552 & 历史单次记录\\
SAR-ShipDataset & Baseline & 94.616 & 49.568 & 非严格配对\\
SAR-ShipDataset & Ours & 94.844 & 57.818 & 非严格配对\\
SADD scratch, 200e & Baseline & 98.335 & 55.018 & 历史单次记录\\
SADD scratch, 200e & P3 & 97.760 & 54.597 & 历史单次记录\\
\bottomrule
\end{tabular}
\end{table}
SAR-AIRcraft test 记录的 AP$_{50:95}$ 差为 $+1.039$ 个百分点，SADD 的差为 $-0.421$ 个百分点。SAR-ShipDataset 的两行不是严格同协议配对，且标注的 AP$_{50:95}$ 差异较大，投稿前必须回到原始 run、评估脚本和配置复核。当前没有完成 SAR-ShipDataset 与 SAR-AIRcraft1.0 的公平三种子矩阵，也没有完成与 YOLOv8n、YOLOv5s 等外部方法的同协议重实现。

\section{讨论}
\subsection{结果回答了什么}
在 SSDD 当前受控设置下，LRC-only 的三种子 AP$_{50:95}$ 描述性均值最高，LSGM-only 的召回均值最高；两者均存在负向 seed。该证据支持把局部响应修正与结构引导作为 SAR 状态建模的设计方向，并明确了组合假设与空间粒度假设的当前边界。full 未超过 LRC-only，且 token-wise 未超过 global/image-wise，是对原始组合和空间粒度假设的直接限制。

跨任务记录显示明显异质性：Farmland 变化检测为正，Yellow River、Sulzberger、两个 scratch 分类配对和两项分割的主 mIoU 记录为负；SAR-ACD 的 VMamba 初始化轨道虽有较高 warmup 结果，但训练日程并不匹配。由于任务头、初始化和实验设置不同，正负方向不能被解释为同一效应量；当前材料尚不足以证明“通用 SAR 视觉提升”。

\subsection{定性可视化与反例}
图\ref{fig:lsgmvis}展示历史 SSDD 验证集审计中的两个固定样例。在一个样例中，GT 框内 importance/background importance 比值为 3.67；另一个样例为 0.45，目标区域 importance 低于背景。完整六样例的均值约 1.19、中位数约 0.72，显示均值受单个高比值样例影响。LSGM 图因此不能称为目标掩膜或稳定注意力解释。
\begin{figure}[H]\centering
\includegraphics[width=.96\linewidth]{figures/fig9_visualization_limits.pdf}
\caption{历史检查点上的引导信息图正向示例与反例。六张验证图按目标数分层固定；此处仅展示分布范围，baseline 与 LSGM 检查点训练协议不匹配，不能作为受控方法证据。}
\label{fig:lsgmvis}
\end{figure}
该审计使用恢复的历史检查点和 validation 图像，不是当前 clean-scratch 三种子矩阵的机制诊断。要支撑机制解释，仍需在同一固定验证协议下比较 baseline、LRC-only、LSGM-only 与 full，并报告完整样本分布和重要性扰动结果。

\subsection{尚未验证的机制解释}
局部残差和空间重要性可能帮助模型对不同空间区域采用不同的扫描输入或状态步长，但 AP 与 mIoU 等终端指标不能证明这一机制。除图\ref{fig:lsgmvis}所示的历史审计外，当前还缺少新矩阵上的跨 seed 预测可视化、预先定义的目标/背景区域统计、重要性扰动或分支置换实验；还需在控制 block 其他部分不变时单独测试残差修正，并评估一维展平趋势卷积跨越图像行边界的影响。

\subsection{审稿人视角下的主要风险}
\begin{enumerate}
\item \textbf{外部比较不足：}除 SSDD 内部变体外，缺少同划分、同 scratch 策略、同训练预算的轻量检测基线重实现；SAR-ShipDataset 与 SAR-AIRcraft1.0 的三种子公平验证未完成。
\item \textbf{结构混杂：}baseline 与 LRC-only 的归一化、MLP ratio 和扫描块结构不完全匹配，当前实验不能隔离局部残差公式的因果作用。
\item \textbf{统计不确定性：}三个 seed 的样本量较小，且 seed 同时改变训练/验证划分；不能声称显著性或普遍稳定性。
\item \textbf{测量口径：}复杂度速度记录缺少完整硬件与计时元数据；提交前需按统一 profiler 补测。
\item \textbf{方法复现细节：}需要冻结代码 commit、插入 stage/stride、参数维度、LSGM 路径、扫描方向和 checkpoint 哈希，并与正文公式逐项对应。
\item \textbf{跨任务证据：}分类/变化检测/分割大多单 seed 或部分可复现，且有负结果；只能用于边界分析，不能当作统一泛化验证。
\end{enumerate}

\subsection{局限与下一步实验}
当前证据最完整的部分是 SSDD 内部消融，不是跨数据集主比较。投稿前优先完成：SAR-ShipDataset 与 SAR-AIRcraft1.0 的 baseline、LRC-only、LSGM-only、full 同协议三种子复验；对 YOLOv8n/YOLOv5s 或同类轻量模型执行公平复现；重新测量完整模型效率；增加控制其余 block 结构不变的局部残差、LSGM、gamma/write、delta 分支消融；复核数据清单与 train/test 场景隔离。若新增实验仍显示 full 低于 LRC-only，应收窄方法 claim 或重新设计组合机制，而不能删去负向结果。

\section{结论}
本文提出了面向 SAR 视觉状态空间模型的 RGSP 方法，在 SSDD 上完成了五种结构变体、三个随机种子的受控消融，并给出了组件单独与组合情形的完整结果：LRC-only 与 LSGM-only 的 AP$_{50:95}$ 描述性均值高于基线，full 未优于 LRC-only，token-wise 引导未优于图像级对照；跨任务诊断结果进一步界定了方法的适用边界。补齐跨数据集公平复验、结构匹配消融、效率协议与机制诊断，是形成最终结论的必要工作。

\appendix
\section{SSDD 逐 seed 原始汇总}
\begin{table}[H]\centering\scriptsize
\caption{SSDD 官方测试集逐 run 指标（\%）。每个单元格对应一个 seed 和一个变体；完整精度及权重 SHA-256 见随稿 CSV。}
\begin{tabular}{llrrrrr}
\toprule
Variant & Seed & AP$_{50}$ & AP$_{75}$ & AP$_{50:95}$ & Precision & Recall\\
\midrule
Baseline & 0 & 87.1822 & 35.6859 & 42.6943 & 89.2285 & 80.4029\\
Baseline & 1 & 83.6934 & 28.1695 & 38.2746 & 87.8971 & 77.2894\\
Baseline & 2 & 88.6714 & 38.2830 & 45.1215 & 91.5192 & 80.0366\\
LRC-only & 0 & 87.1986 & 41.9141 & 45.5944 & 90.9989 & 79.4872\\
LRC-only & 1 & 87.3584 & 35.8419 & 43.5303 & 89.3584 & 79.9723\\
LRC-only & 2 & 85.5283 & 34.9699 & 42.8199 & 89.3093 & 79.5616\\
LSGM-only & 0 & 90.0920 & 38.1567 & 44.8217 & 90.6961 & 83.9133\\
LSGM-only & 1 & 87.3274 & 30.4719 & 41.4304 & 88.6802 & 80.4029\\
LSGM-only & 2 & 87.2713 & 35.8535 & 43.1189 & 91.0193 & 79.8535\\
Full & 0 & 88.0882 & 30.7187 & 42.3400 & 93.2318 & 80.7692\\
Full & 1 & 87.8664 & 36.8629 & 44.2091 & 92.2132 & 78.7546\\
Full & 2 & 87.9535 & 34.7167 & 42.1778 & 90.2943 & 80.0366\\
Global/image-wise & 0 & 86.8336 & 40.0310 & 44.9822 & 91.9194 & 78.0220\\
Global/image-wise & 1 & 86.7447 & 38.0432 & 44.4228 & 90.8392 & 77.8388\\
Global/image-wise & 2 & 85.6446 & 32.5236 & 41.8918 & 92.5226 & 79.4872\\
\bottomrule
\end{tabular}
\end{table}

\section{分割类别 IoU 补充}
FUSAR-Map 的四类 IoU（building / vegetation / water / road）分别为 baseline：14.280 / 71.493 / 86.634 / 4.782\%；P3：15.231 / 71.075 / 86.846 / 1.091\%。road 类下降说明整体 mIoU 之外还存在类别层面的明显退化。

\section{跨任务实验协议与数据覆盖补充}
\subsection{分类}
校正后的 scratch 分类诊断使用输入尺寸 224、batch size 64、最多 200 epochs、seed 0；每个数据集内 baseline 与 P3 使用相同划分。SAR-ACD 的五类协议和 SAMPLE 当前选取的 real-only 子集分别报告，SAMPLE 结果不可与其他采样/合成图像口径混合。另有一条使用 VMamba-Tiny ImageNet 初始化的 SAR-ACD 轨道：骨干使用预训练参数，新加的 P3 模块随机初始化，输入 224、batch 64、最多 200 epochs、seed 0；warmup 20/50 是 P3 专属训练日程。该轨道与 scratch 实验分开，且 warmup 结果不能归因于模块单因素。

\subsection{变化检测}
每个变化检测数据集目前只有一个配准图像对：Ottawa $350\times290$、Yellow River $289\times257$、Farmland $291\times306$、Sulzberger $256\times256$。模型使用密集空间采样以构造训练/验证/测试块，并保持空间区域分离；历史协议说明原始裁块为 $48\times48$ 后缩放到训练输入 $192$。当前汇总记录为 scratch、P3-$d=16$、150 epochs、AdamW/cosine、seed 0，验证集选择模型后在固定测试区域评估。单个场景对无法代表跨场景泛化；同一图像对上的大量裁块也不能视为独立场景重复。

\subsection{语义分割}
FUSAR-Map 数值对应四类 UPerNet-lite 诊断：256 输入块、batch 10、FP32、最多 150 epochs，以验证 mIoU 选模；记录的划分为 17,560/4,390/1,280 个块。该数字来自报告记录，当前稿件尚未附最终 manifest 哈希，不能据此声称已完成场景泄漏审计。该组结果在实验日志中标为部分可复现。OpenEarthMap-SAR 使用 SAR-only 输入、VMamba-Tiny ImageNet 初始化骨干、P3-$d=16$，训练输入 512、测试输入 1024、seed 0，物理 batch 32 加梯度累积 2（等效全局 batch 64），150 epochs；新增模块随机初始化。两项分割结果均保留为边界分析。

\section{单种子与历史检测结果账本}
下表补充主文没有合并的配对检测记录。除表中注明的同协议单 seed 配对外，这些值不构成三种子结论；即使训练设置配对，单次差值也只描述该次运行。
\begin{table}[H]\centering\scriptsize
\caption{单种子或历史配对检测记录（官方/报告测试口径；AP$_{50:95}$ 单位为百分比）。}
\begin{tabularx}{\linewidth}{p{2.2cm}p{2.5cm}rrrX}
\toprule
数据集 & 记录协议 & Baseline & P3/ours & 差值 (pp) & 证据边界\\
\midrule
SSDD & scratch, BS16, 150e & 35.239 & 42.991 & $+7.752$ & seed 0；与 BS32 结果方向相反\\
SSDD & scratch, BS32, 150e & 45.683 & 42.486 & $-3.197$ & seed 0；独立于 BS16\\
SADD & scratch, BS32, 200e & 55.018 & 54.597 & $-0.421$ & seed 0；单次配对\\
AIR-SARShip-2.0 & 源图划分后切片，scratch, BS32, 200e & 17.375 & 19.291 & $+1.916$ & seed 0；文献常用预训练设置不同\\
OGSOD-1.0 & 1024, 300e, global BS64, FP32 & 42.302 & 43.602 & $+1.300$ & 单次匹配 TEST；不支持稳定性结论\\
\bottomrule
\end{tabularx}
\end{table}
\begin{figure}[H]\centering
\includegraphics[width=.9\linewidth]{figures/fig10_historical_detection.pdf}
\caption{单种子历史检测配对差的方向和大小。不同数据集与训练批次分别展示，不计算合并平均值。}
\label{fig:historical-detection}
\end{figure}

以下历史或不完整记录单独列出，不与上述配对表混排：
\begin{table}[H]\centering\scriptsize
\caption{协议不匹配、无配对基线或未完成的检测记录。验证过程值不视为最终测试成绩。}
\begin{tabularx}{\linewidth}{p{2.3cm}p{2.1cm}p{2.6cm}p{2.7cm}X}
\toprule
数据集/协议 & Baseline AP$_{50:95}$ & P3/ours AP$_{50:95}$ & 状态 & 限制\\
\midrule
SSDD 历史 ex27 官方 TEST & 70.348 & 71.445 & 已复核，历史检查点 & 旧训练配方、旧拓扑和初始化入口；不是当前 clean-scratch 对照。历史 VAL 74.022 不得写成 TEST。\\
SSDD 部分迁移 300e & 无匹配项 & 72.261 & 单个 P3 TEST & 不能计算方法增益；不得与当前 150e scratch 矩阵比较。\\
SADD 部分迁移 300e & 无匹配项 & 72.113 & 单个 P3 TEST & 不能计算方法增益；最终验证损失字段另有审计事项。\\
SRSDD-v1.0 HBB & 无匹配项 & 22.147 & 单个 P3 TEST & HBB 转换与公开 OBB 协议不同，缺少匹配 baseline。\\
SAR-ShipDataset & 49.568 & 57.818 & 历史测试记录 & baseline 训练期间 batch 从 256 改为 128，不能视为严格受控配对。\\
SAR-AIRcraft patch TEST & 54.513 & 55.552 & 历史测试记录 & baseline 为 SGD/BS16/约 137e；P3 为 AdamW/全局 BS48/300e；重叠 patch 协议也不同于原图基准。\\
SAR-AIRcraft1.0 失败运行 & -- & -- & 失败，不报告 TEST & FP32 P3 运行在 epoch 15、batch 166 出现非有限 loss；epoch 14 的验证 AP$_{50:95}=53.902$ 只是中间记录，未生成最终测试文件。\\
\bottomrule
\end{tabularx}
\end{table}

SAR-AIRcraft patch 历史记录使用 512 像素重叠裁片。训练、验证和测试分别为 13,976、1,923、2,989 个 patch，模型输入 1024；其 baseline 与 P3 的优化器、batch 和训练长度不匹配。

SAR-ShipDataset 历史记录使用 31,783、3,972、3,974 张训练、验证和测试图像，输入 512；baseline 训练 batch 曾发生变化。以上数字只补充旧运行账本，不替代尚未完成的相同划分、scratch、相同预算三种子复验。

\subsection{已发表检测结果的数值背景（非配对比较）}
为交代当前结果与已发表 SAR 检测工作的距离，表\ref{tab:literature_context}摘录 MaDiNet 作者稿的 SSDD、SAR-ShipDataset 和 SAR-AIRcraft1.0 对照表中的代表性 AP$_{50}$ 数值\cite{madinet_preprint}。这些数值只描述该论文内部的比较；本项目没有在相同划分、初始化、输入、训练预算和评估脚本下重实现这些方法。特别是，SSDD 当前三种子 scratch 结果与 MaDiNet 表中的 AP$_{50}$ 差异较大，不能将其解释为模型优劣的受控估计。SAR-ShipDataset 和 SAR-AIRcraft1.0 的本地值也来自协议不匹配的历史记录。该表用于呈现文献背景和现有差距，不用于主表排名或摘要中的性能优势主张。
\begin{table}[H]\centering\scriptsize
\caption{MaDiNet 作者稿报告的代表性 AP$_{50}$ 与本项目现有证据。跨论文及历史值未做协议匹配，不能直接排名。}
\label{tab:literature_context}
\begin{tabularx}{\linewidth}{p{1.65cm}XXp{3.45cm}}
\toprule
数据集 & MaDiNet 作者稿内报告值（AP$_{50}$，\%） & 本项目现有记录（AP$_{50}$，\%） & 不可直接比较的原因\\
\midrule
SSDD & YOLOv8n 98.1；YOLOv5s 98.3；DiffDet4SAR 96.9；MaDiNet 99.0 & 当前 scratch 矩阵：baseline $86.52\pm2.56$；LRC-only $86.70\pm1.01$；LSGM-only $88.23\pm1.61$；full $87.97\pm0.11$；global $86.41\pm0.66$ & MaDiNet 作者稿与本项目在初始化、切分、预算及评估配置上的一致性未经复核；本项目没有重实现这些行。\\
SAR-ShipDataset & YOLOv5n 89.3；DiffDet4SAR 95.1；MaDiNet 97.5 & 历史 baseline 94.616、P3 94.844；非严格配对 & 本地 batch 变化，且划分、训练设置和评估链未与作者稿对齐。\\
SAR-AIRcraft1.0 & YOLOv8n 88.2；DiffDet4SAR 88.4；MaDiNet 90.5 & 历史 patch 验证：baseline 87.751、P3 88.825；patch 测试：70.402/70.822 & 本地结果使用重叠 patch；baseline 与 P3 的优化器、batch、轮数不一致。\\
\bottomrule
\end{tabularx}
\end{table}
作者稿中的主数据集表及数值定义见其表 3--5；最终投稿若保留此背景表，仍需核对正式出版版本中的相同条目，并补齐本项目的同协议外部基线。

\section{实验完成度与论文主张对应关系}
\begin{table}[H]\centering\scriptsize
\caption{当前证据的完成状态及可支持的论文结论范围。}
\begin{tabularx}{\linewidth}{p{3.3cm}p{2.3cm}X}
\toprule
实验项 & 当前状态 & 可支持范围 / 仍缺的证据\\
\midrule
SSDD 内部消融 & 完成，5 变体 $\times$ 3 seeds & 可报告当前设置下的描述性均值、配对 seed 差和负交互；不足以声称统计显著。\\
SAR-ShipDataset / SAR-AIRcraft1.0 & 正式公平矩阵未完成 & 缺少 baseline、LRC-only、LSGM-only、full 的统一协议三种子结果。\\
外部轻量检测器对比 & 同协议复现未完成 & 文献数只作背景；不能支撑超过 YOLOv8n/YOLOv5s 等主流模型的结论。\\
固定去斑预处理对照 & Lee 等固定滤波对照未完成 & 当前不能判断所提状态更新相对固定去斑预处理的增益。\\
分类、变化检测、分割 & 多为 seed 0 诊断 & 可记录任务依赖的正负方向；不足以主张跨任务普遍泛化。\\
目标尺度/背景复杂度分层 & 尚无预注册分组分析 & 当前不作小目标或复杂背景上的差异性结论。\\
LSGM 的 $\gamma$/write/$\Delta$ 分支机制 & 分支单独消融未完成 & 目前终端指标与历史热图不足以证明具体状态路径的因果机制。\\
模型效率 & 参数量/GFLOPs 有记录，计时元数据不完整 & 不报告可靠速度优劣；须统一硬件、输入、batch、profiler 和重复次数补测。\\
数据清单及场景泄漏核验 & 部分 manifest 可用，最终哈希审计未齐 & 不对全项目数据隔离作已完成声明。\\
\bottomrule
\end{tabularx}
\end{table}

\section{复现材料清单}
当前实验账本要求随论文归档：训练配置与命令、每个 seed 的 split manifest、数据版本标识、源码 commit、逐 run 日志、最终权重 SHA-256、评估命令和原始输出。主 SSDD 结果表和逐 run CSV 已记录权重哈希；正文投稿版仍应给出可审计的匿名归档或可访问的补充材料链接。本文不包含未经核验的数据场景泄露审计结论。
{\footnotesize
\setlength{\itemsep}{0pt}
\begin{thebibliography}{99}
\bibitem{resnet2015}
K. He, X. Zhang, S. Ren, and J. Sun, "Deep Residual Learning for Image Recognition," in \emph{Proc. IEEE Conf. Computer Vision and Pattern Recognition (CVPR)}, 2016, pp. 770--778. arXiv:1512.03385, 2015.
\bibitem{vit2020}
A. Dosovitskiy et al., "An Image is Worth 16x16 Words: Transformers for Image Recognition at Scale," in \emph{Proc. Int. Conf. Learning Representations (ICLR)}, 2021. arXiv:2010.11929, 2020.
\bibitem{vmamba2024}
Y. Liu, Y. Tian, Y. Zhao, H. Yu, L. Xie, Y. Wang, Q. Ye, and Y. Liu, "VMamba: Visual State Space Model," in \emph{Advances in Neural Information Processing Systems (NeurIPS)}, 2024. arXiv:2401.10166.
\bibitem{mamba2023}
A. Gu and T. Dao, "Mamba: Linear-Time Sequence Modeling with Selective State Spaces," arXiv:2312.00752, 2023. \texttt{https://arxiv.org/abs/2312.00752}.
\bibitem{vim2024}
L. Zhu et al., "Vision Mamba: Efficient Visual Representation Learning with Bidirectional State Space Model," arXiv:2401.09417, 2024. \texttt{https://arxiv.org/abs/2401.09417}.
\bibitem{mambayolo2024}
Z. Wang et al., "Mamba YOLO: A Simple Baseline for Object Detection with State Space Model," arXiv:2406.05835, 2024. \texttt{https://arxiv.org/abs/2406.05835}.
\bibitem{rsmamba2024}
Y. Zhao et al., "RS-Mamba for Large Remote Sensing Image Dense Prediction," \emph{IEEE Transactions on Geoscience and Remote Sensing}, 2024. doi:10.1109/TGRS.2024.3425540.
\bibitem{rstmtmamba2025}
Y. Shen et al., "Learning Cross-Task Features With Mamba for Remote Sensing Image Multitask Prediction," \emph{IEEE Transactions on Geoscience and Remote Sensing}, 2025. doi:10.1109/TGRS.2025.3540573.
\bibitem{madinet2025}
J. Zhou et al., "MaDiNet: Mamba Diffusion Network for SAR Target Detection," \emph{IEEE Transactions on Circuits and Systems for Video Technology}, vol. 35, no. 11, pp. 10787--10800, 2025. doi:10.1109/TCSVT.2025.3574657.
\bibitem{madinet_preprint}
J. Zhou et al., "MaDiNet: Mamba Diffusion Network for SAR Target Detection," arXiv:2411.07500v1, 2024, Tables 3--5.
\bibitem{mcmn}
G. Gao, C. Li, X. Zhang, B. Yao, and Z. Chen, "Mamba-CNN hybrid Multi-scale ship detection Network driven by a Dual-perception feature of Doppler and Scattering," \emph{ISPRS Journal of Photogrammetry and Remote Sensing}, vol. 232, pp. 656--674, 2026. doi:10.1016/j.isprsjprs.2026.01.004.
\bibitem{mambadip}
Chinchu A. et al., "MambaDIP: Mamba-Based Deep Image Prior Framework for Despeckling and Enhancement of SAR Images," \emph{IEEE Journal of Selected Topics in Applied Earth Observations and Remote Sensing}, vol. 19, pp. 21255--21271, 2026. doi:10.1109/JSTARS.2026.3702550.
\bibitem{rssnet}
M. Huang, Y. Yang, and Q. Sun, "RSS-Net: A Mamba-Based Network for SAR Image Denoizing," \emph{IEEE Journal of Selected Topics in Applied Earth Observations and Remote Sensing}, vol. 19, pp. 3539--3551, 2026. doi:10.1109/JSTARS.2025.3647971.
\bibitem{specklefree}
S.-W. Chen, X.-C. Cui, X.-S. Wang, and S.-P. Xiao, "Speckle-Free SAR Image Ship Detection," \emph{IEEE Transactions on Image Processing}, vol. 30, pp. 5969--5983, 2021. doi:10.1109/TIP.2021.3089936.
\bibitem{ssdd2017}
J. Li, C. Qu, and J. Shao, "Ship Detection in SAR Images Based on an Improved Faster R-CNN," in \emph{Proc. SAR in Big Data Era}, 2017.
\bibitem{ssdd2021}
T. Zhang et al., "SAR Ship Detection Dataset (SSDD): Official Release and Comprehensive Data Analysis," \emph{Remote Sensing}, vol. 13, no. 18, art. 3690, 2021. doi:10.3390/rs13183690.
\bibitem{sardataset}
Y. Wang et al., "A SAR Dataset of Ship Detection for Deep Learning under Complex Backgrounds," \emph{Remote Sensing}, vol. 11, no. 7, art. 765, 2019. doi:10.3390/rs11070765.
\bibitem{saraircraft}
Z. Wang et al., "SAR-AIRcraft-1.0: High-Resolution SAR Aircraft Detection and Recognition Dataset," \emph{Journal of Radars}, vol. 12, no. 4, pp. 906--922, 2023. doi:10.12000/JR23043.
\bibitem{sarnano}
Y. Zhang et al., "SAR-NanoShipNet: A Scale-Adaptive Network for Robust Small Ship Detection in SAR Imagery," \emph{ISPRS Journal of Photogrammetry and Remote Sensing}, vol. 232, pp. 262--279, 2026. doi:10.1016/j.isprsjprs.2025.12.006.
\bibitem{sfdnet}
B. Wu et al., "A Novel Star Feature Decoupling Network for Multiscale SAR Image Ship Detection," \emph{IEEE Journal of Selected Topics in Applied Earth Observations and Remote Sensing}, vol. 19, pp. 18294--18305, 2026. doi:10.1109/JSTARS.2026.3688771.
\bibitem{fusar}
X. Shi et al., "Object-Level Semantic Segmentation on the High-Resolution Gaofen-3 FUSAR-Map Dataset," \emph{IEEE Journal of Selected Topics in Applied Earth Observations and Remote Sensing}, 2021. doi:10.1109/JSTARS.2021.3063797.
\end{thebibliography}
}
\end{document}
